\section{RELATED WORK}

\paragraph{Robot architectures and orchestration.}
Robot architectures such as CLARAty, CRAM, and hierarchical planning in the now connect planning to execution through runtime feedback and adjustment~\citep{volpe2001claraty,beetz2010cram,kaelbling2011now}. More recent systems organize learned robot capabilities: SayCan grounds skill selection in affordance values~\citep{ichter2023saycan}, Inner Monologue incorporates execution feedback into language planning~\citep{huang2023innermonologue}, and Pigey coordinates subgoals, skill calls, and retries~\citep{galanti2026pigey}. These mechanisms support task-level coordination, but do not directly specify when to interrupt an ongoing policy action or how much physical intervention would improve its continuation.

\paragraph{Test-time control of VLA execution.}
Test-time methods improve deployed VLAs by controlling action selection and execution. V-GPS and RoboMonkey use learned evaluators to select action proposals~\citep{nakamoto2025steering,kwok2025robomonkey}, while CoVer jointly selects instructions and action chunks through alignment scoring~\citep{kwok2026cover}. AutoHorizon adjusts how much of a predicted chunk to execute before replanning~\citep{wang2026autohorizon}. These methods refine execution from the current observation, leaving open how to select a physical intervention when continuing from that state offers poor prospects for success.

\paragraph{Failure detection and recovery.}
Runtime failure detection provides signals for intervention. SAFE predicts failure from VLA features~\citep{gu2025safe}, while FIPER combines observation novelty with action uncertainty~\citep{romer2025fiper}. Such signals alone do not determine where the policy should resume. CycleVLA combines a progress-aware VLA with an external VLM that predicts failures and selects subtasks for backtracking, followed by minimum Bayes risk decoding for retries~\citep{ma2026cyclevla}. Recovery decisions therefore depend on an external VLM's assessment of subtask preconditions, whose connection to the deployed policy's actual ability to recover remains unclear. Moreover, organizing intervention around subtask boundaries makes its timing depend on the chosen task decomposition; these boundaries need not coincide with the moments when intervention would most benefit execution.
